% ===================================================================
% =========================    致谢     =============================
% ===================================================================
\chapter[致谢]{致\quad 谢}\chaptermark{致\quad 谢}% syntax: \chapter[目录]{标题}\chaptermark{页眉}
%\thispagestyle{noheaderstyle}% 如果需要移除当前页的页眉
%\pagestyle{noheaderstyle}% 如果需要移除整章的页眉

此处填写致谢。


\rightline{2026年6月}
\chapter{作者简历及攻读学位期间发表的学术论文与其他相关学术成果}

\section*{作者简历：}
2017年09月——2021年06月，在xxx大学xxx学院xxx系获得学士学位。

2021年09月——2026年06月，在中国科学院xxx研究院攻读博士学位。


\section*{已发表（或正式接受）的学术论文：}

{
\setlist[enumerate]{}% restore default behavior
\begin{enumerate}[nosep]
    \item 论文1，按照参考文献格式 (第一作者，中国科学院分区一区TOP)
    \item 论文2，按照参考文献格式 (第一作者，中国科学院分区一区TOP)
    \item 论文3，按照参考文献格式 (第一作者，CCF A类会议)
\end{enumerate}
}

% \section*{申请或已获得的专利：}

% （无专利时此项不必列出）

\section*{参加的研究项目及获奖情况：}
{
\setlist[enumerate]{}% restore default behavior
\begin{enumerate}[nosep]
    \item \textbf{参与课题一}：xxx
    \item \textbf{参与课题二}：xxx
    \item \textbf{荣誉奖励一}：2024年 xxx奖学金 
    \item \textbf{荣誉奖励二}：2024年 xxx
\end{enumerate}
}


\cleardoublepage[plain]% 让文档总是结束于偶数页，可根据需要设定页眉页脚样式，如 [noheaderstyle]
%---------------------------------------------------------------------------%
